% TEMPLATE-MILC-v3.tex - plantilla maestra para todas las guias MILC v3
% Ver scripts/generadores/build-guia-milc.py para la lista completa de
% claves esperadas. El script valida que no queden placeholders sin
% reemplazar antes de escribir el archivo final.

\documentclass[11pt,letterpaper]{article}
\usepackage[margin=1.75cm]{geometry}
\usepackage[table]{xcolor}
\usepackage{fontspec}
\usepackage[spanish]{babel}
\usepackage{tikz}
\usepackage[most]{tcolorbox}
\usepackage{tabularx}
\usepackage{array}
\usepackage{fancyhdr}
\usepackage{titlesec}
\usepackage{ragged2e}
\usepackage{enumitem}
\usepackage{microtype}

\setmainfont{Helvetica Neue}
\setsansfont{Helvetica Neue}
\setlength{\parindent}{0pt}
\setlength{\parskip}{6pt}
\hyphenpenalty=200
\exhyphenpenalty=200
\pretolerance=400
\tolerance=2000
\setlength{\emergencystretch}{1em}
\renewcommand{\arraystretch}{1.25}
\setlist{leftmargin=5mm,itemsep=1mm,topsep=1mm}
\newcolumntype{Y}{>{\RaggedRight\arraybackslash}X}

\definecolor{milcMagenta}{HTML}{E6007E}
\definecolor{milcVino}{HTML}{5A0038}
\definecolor{milcTurquesa}{HTML}{0093A5}
\definecolor{milcVerde}{HTML}{58A923}
\definecolor{milcMostaza}{HTML}{E5B400}
\definecolor{milcOcre}{HTML}{B86600}
\definecolor{milcNegro}{HTML}{241816}
\definecolor{milcGris}{HTML}{F7F7F4}
\definecolor{milcLinea}{HTML}{E8E2DD}
\definecolor{milcGradePrimary}{HTML}{FF2D87}
\definecolor{milcGradeDark}{HTML}{9F1B56}
\definecolor{milcGradeSoft}{HTML}{FFE0EE}
\definecolor{milcGradeAccent}{HTML}{CFE7FF}
\definecolor{milcGradeText}{HTML}{FFFFFF}
\color{milcNegro}

\pagestyle{fancy}
\fancyhf{}
\lhead{\small Tecnología e Informática · Grado 10}
\rhead{\small Guía 25 · MILC v3}
\cfoot{\small\thepage}
\renewcommand{\headrulewidth}{0pt}

\newtcolorbox{titlebox}[1]{
  enhanced,colback=#1,colframe=#1,arc=7mm,boxrule=0pt,
  left=7mm,right=7mm,top=3mm,bottom=3mm,
  before skip=8pt,after skip=8pt,
  fontupper=\sffamily\bfseries\Large\color{white}
}

\newtcolorbox{softbox}[3]{
  enhanced,colback=#2,colframe=#1,borderline west={4pt}{0pt}{#1},
  arc=2mm,boxrule=.4pt,left=4mm,right=4mm,top=3mm,bottom=3mm,
  before skip=6pt,after skip=8pt,title={#3},coltitle=white,
  colbacktitle=#1,fonttitle=\sffamily\bfseries,fontupper=\RaggedRight,
  attach boxed title to top left={xshift=3mm,yshift=-2mm},
  boxed title style={arc=2mm, boxrule=0pt}
}

\newtcolorbox{drawbox}[2]{
  enhanced,colback=white,colframe=#1,arc=4mm,boxrule=1pt,
  left=5mm,right=5mm,top=4mm,bottom=4mm,before skip=8pt,after skip=8pt,
  title={#2},coltitle=white,colbacktitle=#1,fonttitle=\sffamily\bfseries,
  fontupper=\RaggedRight,
  attach boxed title to top left={xshift=4mm,yshift=-2mm},
  boxed title style={arc=3mm, boxrule=0pt}
}

\newtcolorbox{infoband}[1]{
  enhanced,colback=#1!8,colframe=#1!50,arc=2mm,boxrule=.5pt,
  left=4mm,right=4mm,top=2mm,bottom=2mm,before skip=4pt,after skip=4pt,
  fontupper=\small\RaggedRight
}

% Puente narrativo entre secciones (contrato editorial Conéctate).
% Da continuidad: "Acabas de X. Ahora vas a Y porque Z."
% Visualmente: banda izquierda en color del grado, texto en itálica pequeña.
\newtcolorbox{puentebox}{
  enhanced,colback=milcGradeSoft,colframe=milcGradeDark,
  borderline west={3pt}{0pt}{milcGradeDark},
  arc=0mm,boxrule=0pt,left=5mm,right=4mm,top=2.5mm,bottom=2.5mm,
  before skip=4pt,after skip=8pt,
  fontupper=\itshape\small\color{milcNegro}\RaggedRight
}
\newcommand{\puente}[1]{%
  \begin{puentebox}%
    \textcolor{milcGradeDark}{\textbf{\textrightarrow}}\hspace{2mm}#1%
  \end{puentebox}%
}

\newcommand{\linead}[1]{\noindent\color{milcLinea}\rule{#1}{0.5pt}\color{milcNegro}}
\newcommand{\respuesta}[1]{\par\vspace{2mm}\foreach \n in {1,...,#1}{\noindent\color{milcLinea}\rule{\linewidth}{0.5pt}\par\vspace{5mm}}\color{milcNegro}}
\newcommand{\checkbox}{\textcolor{milcGradeDark}{\fbox{\phantom{x}}}\hspace{5pt}}

\newcommand{\phasecard}[4]{
\begin{tcolorbox}[enhanced,colback=#3,colframe=#2,arc=3mm,boxrule=.5pt,height=3.05cm,valign=center,halign=center,left=1.5mm,right=1.5mm,top=2mm,bottom=2mm]
{\sffamily\bfseries\fontsize{22}{24}\selectfont\textcolor{#2}{#1}}\par
{\sffamily\bfseries\small #4}
\end{tcolorbox}
}

\begin{document}
\thispagestyle{empty}
\begin{tikzpicture}[remember picture,overlay]
  \fill[white] (current page.south west) rectangle (current page.north east);
  \path[fill=milcGradePrimary,rounded corners=16mm]
    ([xshift=.55cm,yshift=-.55cm]current page.north west) rectangle
    ([xshift=-.55cm,yshift=3.65cm]current page.south east);

  \node[anchor=north west,text=milcGradeText] at ([xshift=2.0cm,yshift=-1.85cm]current page.north west)
    {\sffamily\bfseries\fontsize{14}{16}\selectfont GRADO};
  \node[anchor=north west,text=milcGradeText,opacity=.82] at ([xshift=2.0cm,yshift=-2.75cm]current page.north west)
    {\sffamily\bfseries\fontsize{9}{11}\selectfont SERIE GUÍAS · TECNOLOGÍA E INFORMÁTICA};

  \begin{scope}[shift={([xshift=-3.05cm,yshift=-2.55cm]current page.north east)}]
    \fill[white,opacity=.80] (-.62,-.52) rectangle (.62,.52);
    \draw[milcGradeText,line width=1pt] (-.62,-.52) rectangle (.62,.52);
    \fill[milcGradeDark] (-.42,-.38) rectangle (-.22,.06);
    \fill[milcGradeAccent] (-.10,-.38) rectangle (.10,.24);
    \fill[milcGradePrimary!60!black] (.22,-.38) rectangle (.42,.38);
  \end{scope}

  \node[anchor=west,text=milcGradeText] at ([xshift=1.9cm,yshift=-12.85cm]current page.north west)
    {\sffamily\bfseries\fontsize{124}{124}\selectfont 10};
  \draw[milcGradeText,line width=1.7pt]
    ([xshift=4.52cm,yshift=-11.68cm]current page.north west) circle (.13cm);

  \node[anchor=west,text=milcGradeText,text width=8.7cm,align=left] at ([xshift=10.35cm,yshift=-11.6cm]current page.north west)
    {
     {\sffamily\bfseries\fontsize{19}{22}\selectfont Visualización ---\\gráficos contables\\honestos con IA}\\[4mm]
     {\sffamily\bfseries\fontsize{23}{26}\selectfont Guía 25-10-TIC}\\[2mm]
     {\sffamily\fontsize{13}{16}\selectfont Período 3 · Ofimática con IA: contabilidad, datos y emprendimiento}\\[5mm]
     {\sffamily\bfseries\fontsize{11}{13}\selectfont Institución Educativa Sor María Juliana}\\[1mm]
     {\sffamily\fontsize{10}{12}\selectfont Autor: Dr. Álvaro Cárdenas Orozco}
    };

  \path[fill=milcGradeSoft,rounded corners=7mm]
    ([xshift=1.35cm,yshift=.85cm]current page.south west) rectangle
    ([xshift=-1.35cm,yshift=4.25cm]current page.south east);
  \draw[milcGradeDark,line width=.65pt,rounded corners=7mm]
    ([xshift=1.35cm,yshift=.85cm]current page.south west) rectangle
    ([xshift=-1.35cm,yshift=4.25cm]current page.south east);
  \node[anchor=center,text=milcNegro,text width=17.0cm,align=center] at ([yshift=2.55cm]current page.south)
    {\sffamily\fontsize{10}{12}\selectfont Metodología MILC · Apertura ancestral · Pensamiento computacional · Triángulo Dussel-Estoicismo-Floridi\\
    \textcolor{milcGradeDark}{\bfseries Producto:} 3 gráficos hechos sobre tus datos contables reales del periodo (sesiones 2-4) en Google Sheets (a) gráfico de barras comparando egresos por categoría, (b) gráfico de líneas mostrando evolución del balance semanal, (c) gráfico circular con proporción de cada categoría sobre el total + crítica analítica de 1 gráfico engañoso real (publicidad, prensa o redes sociales) identificando al menos 3 errores de honestidad visual + propuesta de cómo lo redibujarías para que sea honesto};
\end{tikzpicture}
\mbox{}
\newpage

\section*{Apertura: Visualización de datos contables con IA --- gráficos honestos}

\begin{softbox}{milcGradeDark}{milcGradeSoft}{Saber ancestral}
En las plazas de mercado del centro de Cartago, en la galería de Pereira, en las ferias del Quindío, hubo una figura que sostenía la honestidad del comercio: \textbf{el pregonero del mercado}. No era el vendedor; era el que pregonaba productos y precios desde una caseta alta. \emph{``Llegó la naranja del Valle a mil quinientos el kilo, llegó el café del Quindío a quince mil la libra''}. Pero hacía algo más: cuando un comprador se acercaba con duda sobre peso o precio, el pregonero \textbf{mostraba la pesa visible al cliente}. Sacaba la balanza, ponía el producto, dejaba que el cliente leyera la cifra. La pesa visible era \textbf{prueba pública de honestidad}: \emph{``vea, dos kilos justos''}, decía señalando el contrapeso. El comprador miraba, aprobaba, pagaba. La sabiduría era inquebrantable: \textbf{el comerciante que ocultaba la pesa perdía la clientela del barrio entero en una semana}. Esa práctica ancestral se traduce al gráfico contable moderno: \textbf{el gráfico honesto muestra los datos como son; el manipulado los esconde o exagera}.
\end{softbox}

\begin{softbox}{milcTurquesa}{milcGris}{Saber contemporáneo}
La \textbf{visualización de datos} es la práctica de convertir números en imágenes que comuniquen información de manera rápida y precisa. Sus 3 tipos básicos para datos contables: (1) \textbf{Gráfico de barras}: ideal para comparar valores entre categorías (cuánto se gastó en cada categoría este mes). (2) \textbf{Gráfico de líneas}: ideal para mostrar evolución en el tiempo (balance semanal a lo largo del periodo). (3) \textbf{Gráfico circular} (pastel): ideal para proporciones del total (qué porcentaje de los egresos fue alimentación). En Google Sheets, los gráficos se generan automáticamente: seleccionar tabla → Insertar → Gráfico → Sheets sugiere el tipo o eliges tú. Los \textbf{5 principios de honestidad visual} son irrenunciables: (a) \textbf{Eje Y desde cero} en gráficos de barras (salvo justificación explícita). (b) \textbf{Escala consistente}: si el eje muestra miles, todos los valores en miles. (c) \textbf{Sin 3D engañoso}: el efecto 3D distorsiona proporciones. (d) \textbf{Etiquetas legibles}: cada barra o segmento con su nombre y valor. (e) \textbf{Título descriptivo}: que diga qué muestra el gráfico, no solo \emph{``Gráfico 1''}. La regla profesional: \textbf{si tu gráfico necesita explicación verbal extensa, está mal diseñado}.
\end{softbox}

\begin{softbox}{milcMostaza}{milcGris}{Pregunta-puente}
\textit{¿Qué sabía el pregonero al mostrar la pesa visible al comprador, que el novato olvida cuando publica gráficos con eje Y truncado? ¿Y por qué los 5 principios de honestidad visual valen más que las plantillas elegantes mal aplicadas?}
\end{softbox}

\begin{softbox}{milcVerde}{milcGris}{Saber y saber hacer}
Al terminar podrás: (1) \textbf{identificar} cuándo conviene barras, líneas o circular según el tipo de dato y la pregunta; (2) \textbf{analizar} un gráfico real (publicitario o periodístico) detectando errores de honestidad visual; (3) \textbf{crear} 3 gráficos honestos en Sheets sobre tus datos contables reales; (4) \textbf{evaluar} tu propio gráfico aplicando los 5 principios.
\end{softbox}

\small
\begin{tabularx}{\textwidth}{>{\bfseries}p{3.0cm}Y}
\rowcolor{milcGradePrimary!12}Autor & Dr. Álvaro Cárdenas Orozco\\
DBA & Aplicación de ofimática asistida por IA a contabilidad básica y emprendimiento (MEN, T\&I)\\
Referentes & Lean Startup · Phronesis financiera · Ética del dato empresarial\\
Producto final & 3 gráficos hechos sobre tus datos contables reales del periodo (sesiones 2-4) en Google Sheets (a) gráfico de barras comparando egresos por categoría, (b) gráfico de líneas mostrando evolución del balance semanal, (c) gráfico circular con proporción de cada categoría sobre el total + crítica analítica de 1 gráfico engañoso real (publicidad, prensa o redes sociales) identificando al menos 3 errores de honestidad visual + propuesta de cómo lo redibujarías para que sea honesto\\
\end{tabularx}
\normalsize

\puente{Antes de empezar, mira el viaje completo. En las sesiones 2-4 produjiste contabilidad personal y proyección financiera. Hoy aprendes a \textbf{comunicar esos datos visualmente con honestidad}. Las próximas sesiones (encuestas de mercado, Canvas, comunicación empresarial, proyecto) usarán los gráficos honestos como base de presentación.}

\begin{titlebox}{milcGradeDark}
Ruta MILC: así vamos a aprender
\end{titlebox}
\begin{tabularx}{\textwidth}{>{\centering\arraybackslash}X>{\centering\arraybackslash}X>{\centering\arraybackslash}X>{\centering\arraybackslash}X}
\phasecard{1}{milcVerde}{milcGris}{Escuta\\[-1mm]\small Observo, escucho y pregunto.} &
\phasecard{2}{milcTurquesa}{milcGris}{Sistematización\\[-1mm]\small Pensamiento computacional.} &
\phasecard{3}{milcMagenta}{milcGris}{Praxis\\[-1mm]\small Construyo y aplico.} &
\phasecard{4}{milcVino}{milcGris}{Evaluación\\[-1mm]\small Reflexión liberadora.}
\end{tabularx}


\puente{Antes de teorizar, vas a hacer una \emph{auditoría rápida}: busca 1 gráfico real en una publicidad reciente (puede ser de un banco, una tienda, una marca de comida) y otro 1 en una noticia de periódico o canal. Anota qué te llama la atención de cada uno.}

\begin{titlebox}{milcVerde}
Fase 1 · Escuta
\end{titlebox}

\begin{softbox}{milcVerde}{milcGris}{Escena para escuchar}
\textbf{Actividad 1 · IDENTIFICA --- Auditoría de 2 gráficos reales} (15 min · individual con dispositivo). Busca en internet o redes: (1) \textbf{Gráfico publicitario}: una publicidad de banco, tienda, marca, política. (2) \textbf{Gráfico periodístico}: una noticia de prensa o canal con gráfico. Para cada uno, anota: (a) ¿qué tipo de gráfico es (barras, líneas, circular, otro)? (b) ¿el eje Y empieza en cero o en otro valor? (c) ¿hay efecto 3D? (d) ¿qué te llama la atención? (e) ¿alguno parece exagerar o esconder algo?.
\end{softbox}

\checkbox Encontré 2 gráficos reales (publicitario + periodístico).\\
\checkbox Anoté tipo, ejes, efectos visuales.\\
\checkbox Detecté si alguno parece manipulado.

\begin{infoband}{milcVerde}
\textbf{Lo que va al cuaderno (Actividad 1 · IDENTIFICA).} Título: «Actividad 1 --- 2 gráficos bajo la lupa». \textbf{Formato:} captura o referencia de cada gráfico + 5 preguntas respondidas. \textbf{Sabes que terminaste cuando} ves que la honestidad visual no se da por sentada.
\end{infoband}

\puente{Ya tienes 2 gráficos reales. Ahora vas a entender la \textbf{anatomía de los 3 tipos básicos} y los 5 principios de honestidad visual.}

\begin{titlebox}{milcTurquesa}
Fase 2 · Sistematización con pensamiento computacional
\end{titlebox}

\begin{softbox}{milcTurquesa}{milcGris}{Cuatro pilares aplicados al tema}
Tu auditoría te muestra que la honestidad varía. Ahora necesitas la \textbf{anatomía profesional} de los 3 tipos básicos y los 5 principios de honestidad visual.
\end{softbox}

\small
\begin{tabularx}{\textwidth}{>{\bfseries\RaggedRight\arraybackslash}p{3.6cm}Y}
\rowcolor{milcTurquesa!90}\color{white}Pilar & \color{white}\bfseries Aplicación al tema\\
1. Descomposición & Los \textbf{3 tipos básicos} se descomponen así: (1) \textbf{Barras}: comparar valores entre categorías. Para tus egresos: cuánto en transporte vs alimentación vs ocio. Las barras pueden ser verticales u horizontales. (2) \textbf{Líneas}: mostrar evolución en el tiempo. Para tu balance: cómo cambió de la semana 1 a la 2 a la 3. Las líneas conectan puntos en orden cronológico. (3) \textbf{Circular} (pastel): mostrar proporción del total. Para tus egresos: qué porcentaje del total fue alimentación, transporte, ocio. Solo funciona con pocas categorías (3-6 máximo) y cuando la suma tiene sentido conceptual.\\
2. Reconocimiento de patrones & Los \textbf{5 principios de honestidad visual} son ley del oficio: (a) \textbf{Eje Y desde cero}: en gráficos de barras siempre. Empezar el eje en valor no cero exagera diferencias pequeñas. (b) \textbf{Escala consistente}: si miden miles, todos en miles. No mezclar millones con miles sin etiquetar. (c) \textbf{Sin 3D engañoso}: el efecto 3D distorsiona proporciones reales. Mantenerse en 2D. (d) \textbf{Etiquetas legibles}: nombres de categorías y valores numéricos visibles. (e) \textbf{Título descriptivo}: que diga qué muestra el gráfico (\emph{``Egresos por categoría, semana 1-4''} mejor que \emph{``Mi gráfico''}).\\
3. Abstracción & El patrón profesional fundamental: \textbf{cada tipo de gráfico responde una pregunta distinta}. Barras responden \emph{``¿cuánto es cada uno?''}. Líneas responden \emph{``¿cómo cambió en el tiempo?''}. Circular responde \emph{``¿qué porcentaje del total?''}. Elegir mal el tipo distorsiona la comunicación. La phronesis del oficio consiste en \textbf{elegir el tipo de gráfico según la pregunta, no según la estética}.\\
4. Algoritmo conceptual & Algoritmo: (1) en tu plantilla de Sheets, seleccionar la tabla de datos. (2) Insertar → Gráfico. (3) Sheets sugiere tipo; verificar que es el correcto para tu pregunta. (4) verificar eje Y desde cero. (5) editar título descriptivo. (6) agregar etiquetas si Sheets no lo hizo. (7) eliminar efectos 3D si los hay. (8) verificar legibilidad imprimiendo o viendo en celular.\\
\end{tabularx}
\normalsize

\begin{drawbox}{milcTurquesa}{Anatomía de gráficos honestos}
\textbf{Barras:} comparar valores entre categorías. \\[1mm] \textbf{Líneas:} mostrar evolución en el tiempo. \\[1mm] \textbf{Circular:} proporciones del total (máx 3-6 categorías). \\[1mm] \textbf{5 principios:} eje cero / escala / no 3D / etiquetas / título.
\end{drawbox}

\begin{softbox}{milcMostaza}{milcGris}{Errores comunes y su corrección}
(1) \textbf{Eje Y truncado}: empezar en valor distinto de cero exagera. (2) \textbf{3D que distorsiona}: las proporciones cambian con el ángulo. (3) \textbf{Circular con muchas categorías}: ilegible con más de 6. (4) \textbf{Tipo equivocado}: usar circular para evolución temporal. (5) \textbf{Sin título o título vago}: imposible saber qué muestra. (6) \textbf{Sin fuente declarada}: imposible verificar los datos. (7) \textbf{Colores arbitrarios sin propósito}.
\end{softbox}

\begin{infoband}{milcTurquesa}
\textbf{Lo que va al cuaderno (Actividad 2 · ANALIZA).} Título: «Actividad 2 --- Anatomía gráficos honestos». \textbf{Formato:} (a) 3 tipos con uso típico; (b) 5 principios; (c) los 7 errores con marca del que más temes. \textbf{Sabes que terminaste cuando} podrías auditar cualquier gráfico real.
\end{infoband}

\puente{Tienes el método. Toca aplicarlo: generas 3 gráficos honestos en Sheets con tus datos contables, criticas 1 gráfico engañoso real, propones cómo redibujarlo.}

\begin{titlebox}{milcMagenta}
Fase 3 · Praxis con pensamiento computacional
\end{titlebox}

\begin{softbox}{milcMagenta}{milcGris}{Construyendo el producto con los cuatro pilares}
\textbf{Actividad 3 · CREA + EVALÚA --- 3 gráficos honestos + crítica + propuesta} (40 min · individual con dispositivo). Hora de aplicar. Generas 3 gráficos en Sheets con tus datos contables, criticas 1 gráfico engañoso real, propones cómo redibujarlo.
\end{softbox}

\small
\begin{tabularx}{\textwidth}{>{\bfseries\RaggedRight\arraybackslash}p{3.6cm}Y}
\rowcolor{milcMagenta!90}\color{white}Pilar & \color{white}\bfseries Aplicación al producto\\
Descomposición & Tus 3 gráficos se construyen en 5 pasos: (1) abrir tu plantilla Sheets contable con datos de las 2 semanas. (2) \textbf{Gráfico de barras}: seleccionar columna de categorías + columna de totales por categoría, Insertar Gráfico. (3) \textbf{Gráfico de líneas}: seleccionar fechas + columna de balance acumulado por semana. (4) \textbf{Gráfico circular}: seleccionar categorías + totales (Sheets sugiere circular si las categorías son pocas). (5) editar títulos descriptivos y verificar los 5 principios de honestidad en los 3.\\
Patrones & La \textbf{crítica de 1 gráfico engañoso} se hace tomando un gráfico real (publicitario o periodístico) e identificando concretamente: (a) qué principio de honestidad viola. (b) qué impresión falsa genera al lector. (c) cómo lo redibujarías para que sea honesto. Ejemplo: \emph{``La publicidad del banco X muestra que las tasas de interés bajaron 0.5 por ciento, pero el gráfico hace parecer que bajaron 50 por ciento al usar eje Y desde 5 por ciento''}.\\
Abstracción & Sigue el patrón \textbf{``honestidad antes que estética''}: un gráfico funcional y honesto vale más que uno bonito y engañoso. La phronesis del oficio consiste en \textbf{rechazar el adorno que oculta la verdad}.\\
Algoritmo de armado & Algoritmo: (1) 3 gráficos en Sheets con tus datos. (2) verificación de los 5 principios en cada uno. (3) elección de gráfico engañoso real. (4) análisis de errores. (5) propuesta de redibujado.\\
\end{tabularx}
\normalsize

\begin{drawbox}{milcMagenta}{Checklist antes de entregar}
\checkbox 3 gráficos sobre datos contables reales. \\[1mm] \checkbox Tipo apropiado para cada pregunta. \\[1mm] \checkbox Eje Y desde cero en barras y líneas. \\[1mm] \checkbox Sin efecto 3D. \\[1mm] \checkbox Títulos descriptivos en los 3. \\[1mm] \checkbox 1 gráfico engañoso analizado con 3 errores. \\[1mm] \checkbox Propuesta de redibujado clara.
\end{drawbox}

\begin{softbox}{milcMagenta}{milcGris}{Plantilla de trabajo}
\textbf{Gráfico 1 (barras).} \textit{Egresos por categoría.} \\[1mm] \textbf{Gráfico 2 (líneas).} \textit{Balance semanal.} \\[1mm] \textbf{Gráfico 3 (circular).} \textit{Proporción de egresos.} \\[1mm] \textbf{Crítica.} \textit{Gráfico real + 3 errores + redibujado.}
\end{softbox}

\begin{infoband}{milcMagenta}
\textbf{Lo que va al cuaderno (Actividad 3 · CREA + EVALÚA).} Título: «Actividad 3 --- 3 gráficos honestos + crítica». \textbf{Formato:} (a) capturas de los 3 gráficos; (b) gráfico engañoso analizado; (c) propuesta de redibujado. \textbf{Sabes que terminaste cuando} tus 3 gráficos comunican tus datos con honestidad.
\end{infoband}

\puente{Lo que sigue resume \emph{qué} entregas hoy: 3 gráficos honestos + crítica de 1 engañoso + propuesta de corrección. El criterio principal: que tus gráficos comuniquen los datos con honestidad y que sepas reconocer cuándo otros mienten visualmente.}

\begin{titlebox}{milcMostaza}
Producto de la sesión
\end{titlebox}
\textbf{3 gráficos honestos sobre datos reales + crítica + redibujado}.
\begin{softbox}{milcMostaza}{milcGris}{Criterios de calidad}
(1) 3 tipos distintos (barras, líneas, circular). (2) Datos contables reales del periodo. (3) 5 principios de honestidad aplicados. (4) Crítica de gráfico engañoso con 3 errores. (5) Propuesta clara de redibujado. (6) Títulos descriptivos en todos.
\end{softbox}

\puente{Antes de cerrar, mira los gráficos desde las cinco dimensiones humanas. Mostrar los datos como son es respeto por el lector; manipular visualmente es traición al oficio de la información.}

\begin{titlebox}{milcVino}
Fase 4 · Evaluación liberadora — cinco dimensiones
\end{titlebox}

\begin{softbox}{milcVino}{milcGris}{Sentido de la evaluación}
Evaluar no es solo revisar una nota. Es reconocer qué aprendí, qué cambió en mí, qué emociones manejé, cómo me relaciono con la ciudad y la comunidad, y qué le dejaré a quien viene detrás.
\end{softbox}

\textbf{Mi reflexión liberadora en cinco dimensiones.} Completa cada frase iniciada:

\small
\begin{tabularx}{\textwidth}{>{\bfseries\RaggedRight\arraybackslash}p{3.4cm}Y>{\RaggedRight\arraybackslash}p{4.6cm}}
\rowcolor{milcVino}\color{white}Dimensión & \color{white}\bfseries Mi respuesta (frame starter) & \color{white}\bfseries Algo que descubrí\\
1. Personal & Lo que más cambió en mí fue \linead{6.5cm} & \linead{4.2cm}\\
2. Emocional & La emoción más fuerte fue \linead{2.5cm} y la regulé \linead{3cm} & \linead{4.2cm}\\
3. Ciudadana & En democracia esto importa porque \linead{6cm} & \linead{4.2cm}\\
4. Local (Cartago) & En mi barrio sería útil para \linead{6.5cm} & \linead{4.2cm}\\
5. Intergeneracional & A mi abuela/abuelo le diría \linead{6.5cm} & \linead{4.2cm}\\
\end{tabularx}
\normalsize

\begin{infoband}{milcVino}
\textbf{Para la próxima sustentación.} Vuelve a esta tabla en seis meses. Verás que las respuestas cambian — no porque lo aprendido cambie, sino porque tú cambias. La evaluación liberadora es una conversación contigo a través del tiempo.
\end{infoband}


\puente{Y para terminar, tres voces sobre la honestidad visual. Dussel sobre los gráficos que respetan al lector apurado, Marco Aurelio sobre la simplicidad como virtud visual, Floridi sobre la visualización honesta como ética informacional.}

\begin{titlebox}{milcGradeDark}
Triángulo de pensamiento · cierre filosófico
\end{titlebox}

\begin{softbox}{milcVino}{milcGris}{Dussel · lente del nosotros}
\textit{``Un gráfico honesto respeta al lector apurado; uno manipulado le quita el tiempo y le miente.''}

La mayoría de las personas que verán tus gráficos los miran en redes, prensa rápida, presentaciones con luz baja. La filosofía de la liberación pide diseñar para esa condición: \textbf{gráficos que comuniquen verdad al primer vistazo}. La phronesis del oficio visual consiste en \textbf{no engañar a quien no tiene tiempo de verificar}.

\textbf{¿Mis gráficos respetan al lector apurado con verdad clara, o lo engañan exagerando con trucos visuales?}
\end{softbox}
\linead{\linewidth}\\[1mm]
\linead{\linewidth}

\begin{softbox}{milcMostaza}{milcGris}{Estoicismo · Marco Aurelio · cuidado interior}
\textit{``Mostrar los datos como son es virtud visual; manipular para impactar es vicio del comunicador.''}

Es tentador exagerar diferencias para que el gráfico \emph{``impresione''}. La sabiduría estoica recomienda lo opuesto: \emph{muestra las diferencias como son, ni más grandes ni más pequeñas}. Si la diferencia es del 2 por ciento, el gráfico debe mostrarla del 2 por ciento. El impacto visual no debe contradecir la verdad de los datos.

\textbf{¿Mis gráficos muestran los datos como son, o los exagero para que impresionen?}
\end{softbox}
\linead{\linewidth}\\[1mm]
\linead{\linewidth}

\begin{softbox}{milcTurquesa}{milcGris}{Floridi · lente de la infoesfera}
\textit{``La visualización honesta es ética informacional en la era de los gráficos virales.''}

En la era de la información, los gráficos circulan más rápido que los datos. La gente comparte sin verificar. La ética digital pide que quien produce gráficos respete esa velocidad de circulación produciendo gráficos honestos que sigan siendo correctos al compartirse sin contexto.

\textbf{Si mis gráficos se compartieran en redes sin mi explicación, ¿seguirían comunicando la verdad?}
\end{softbox}
\linead{\linewidth}\\[1mm]
\linead{\linewidth}

\begin{titlebox}{milcVino}
Compromiso final
\end{titlebox}

\small
\begin{tabularx}{\textwidth}{>{\RaggedRight\arraybackslash}p{3.0cm}YYY}
\rowcolor{milcVino}\color{white}\bfseries Criterio & \color{white}\bfseries Lo logré & \color{white}\bfseries En proceso & \color{white}\bfseries Necesito apoyo\\
Comprensión del tema & Explico ideas centrales con mis palabras. & Reconozco ideas pero falta precisión. & Necesito repasar conceptos base.\\
Pensamiento computacional & Apliqué los 4 pilares al diseñar mi producto. & Apliqué algunos pilares. & Necesito releer la fase 2.\\
Aplicación y producto & Mi producto cumple los criterios y se puede revisar. & Producto funcional con ajustes pendientes. & Aún en construcción.\\
Reflexión 5 dimensiones & Respondí cada dimensión con honestidad. & Respondí algunas con detalle. & Mis respuestas son muy generales.\\
Triángulo de pensamiento & Conecté las 3 voces con mi trabajo real. & Conecté al menos una. & No relaciono las voces con mi proceso.\\
\end{tabularx}
\normalsize

\textbf{Hoy entendí que:}\\[1mm]
\linead{\linewidth}\\[1mm]
\linead{\linewidth}

\textbf{Mi compromiso para la próxima sesión es:}\\[1mm]
\linead{\linewidth}\\[1mm]
\linead{\linewidth}

\begin{softbox}{milcMostaza}{milcGris}{Cita de despedida · Marco Aurelio}
\textit{``El obstáculo en el camino se vuelve el camino. Lo que estorba al camino se vuelve camino.''}\\[1mm]
Cada error de hoy, bien observado, se convierte en pieza del camino que pisarás mañana.
\end{softbox}

\small
\begin{tabularx}{\textwidth}{>{\bfseries}p{3.6cm}Y}
\rowcolor{milcGradeDark!12}Nombre completo: & \linead{10cm}\\
Fecha: & \linead{4cm} \quad Curso/grupo: \linead{4cm}\\
Mi firma: & \linead{9cm}\\
\end{tabularx}
\normalsize

\begin{infoband}{milcGradeDark}
\textbf{Para la próxima sesión.} Llega con tus 3 gráficos honestos y la crítica del engañoso. En la sesión 6 vas a aprender a hacer \textbf{encuestas y análisis de mercado con IA} para microemprendimiento. Los gráficos contables de hoy son tu economía personal; las encuestas de mercado de mañana son el ojo a tus clientes potenciales.
\end{infoband}

\end{document}
